% Lösungen Einheit 3 von 4 – Lotfußpunkt (zusammenbau v0.1)
\einheitenkopf{Einheit 3 von 4 · Lotfußpunkt}

\erg{14}{a) Lotbedingung (Skalarprodukt null) \quad b) $F_t(2 + 2t | 1 - t | 2t)$, $\overrightarrow{PF_t} \circ (2 | -1 | 2) = 9t - 9 = 0$, $t = 1$: F(4 | 0 | 2); Abstand $|\overrightarrow{PF}| = |(0 | -2 | -1)| \approx 2{,}24$ \quad c) Q ist der Lotfußpunkt des Lots von E auf die Gerade durch A und B: I legt Q auf diese Gerade, II macht $\overrightarrow{EQ}$ senkrecht zu ihr. \quad d) M als Mittelpunkt von AB berechnen; R als allgemeinen Punkt $R_t$ der Geraden CD ansetzen; das Skalarprodukt von $\overrightarrow{MR_t}$ und $\overrightarrow{CD}$ null setzen und nach t lösen; t einsetzen – R ist der Lotfußpunkt, nicht der Mittelpunkt von CD. \quad e) $P_t(6 - 4t | 1 + 4t | 3t)$ auf BE; $\overrightarrow{P_tA} \circ \overrightarrow{BE} = 16 - 41t = 0$, $t = \tfrac{16}{41}$; $|\overrightarrow{PA}| \approx 3{,}12$; wegen der Symmetrie ist $|\overrightarrow{PC}| = |\overrightarrow{PA}|$, Weg $\approx 6{,}25$ cm \quad f) Lotgerade durch (3 | 3 | 3) mit Richtung (1 | 2 | 2) schneidet E für $t = -1$ in F(2 | 1 | 1); gesuchte Gerade: $\vec x = (2 | 1 | 1) + s \cdot (2 | -1 | 0)$ \quad g) Skizze: rechtwinkliges Dreieck mit der Kathete 0,3 (Abstand zu k) und der Hypotenuse QS: $|QS| = \tfrac{0{,}3}{\mathrm{sin}\,30^\circ} = 0{,}6$; $S = Q - 0{,}6 \cdot \tfrac{1}{5}(3 | 4 | 0) = (5{,}64 | 7{,}52 | 0)$ \quad h) „Das Dreieck wird um die Gerade AB gedreht, bis C senkrecht über der xy-Ebene liegt. Bestimme den Bildpunkt T von C.“ S(3 | 1 | 0) ist der Lotfußpunkt von C auf AB, der Mittelpunkt des Drehkreises.}
\erg{15}{a) Fehler: Mittelpunkt statt Lotfußpunkt genommen. Richtig: $R_t(4t | 1 | 6 - 4t)$, $\overrightarrow{MR_t} \circ (4 | 0 | -4) = 32t - 20 = 0$, $t = 0{,}625$, R(2,5 | 1 | 3,5) \quad b) Eingesetzt entsteht eine lineare Gleichung in t mit dem Faktor $|\vec r|^2 \ne 0$ vor t; sie hat genau eine Lösung, und der zugehörige Punkt liegt auf der Geraden mit $\overrightarrow{PF}$ senkrecht zu ihr. \quad c) Lotfußpunkt auf der Achse L(2 | 3 | 1,2), Radius $|\overrightarrow{LA}| = 5$ m; Weg je Umdrehung $2\pi \cdot 5 \approx 31{,}4$ m}
